\documentclass[11pt,a4paper,fontset=fandol]{ctexart}

\usepackage[a4paper,top=2.25cm,bottom=2.45cm,left=2.25cm,right=2.25cm,headheight=18pt,headsep=0.55cm,footskip=24pt]{geometry}
\usepackage{amsmath,amssymb,amsthm}
\usepackage{fancyhdr}
\usepackage{lastpage}
\usepackage{enumitem}
\usepackage{titlesec}
\usepackage{xcolor}
\usepackage{tikz}
\usepackage{hyperref}
\usepackage{bookmark}
\usepackage[most]{tcolorbox}
\usepackage{setspace}
\usepackage{array}
\usetikzlibrary{arrows.meta}

\hypersetup{
  colorlinks=true,
  linkcolor=blue!50!black,
  urlcolor=blue!50!black
}

\setstretch{1.13}
\setlength{\parindent}{2em}
\setlength{\parskip}{0.25em}
\allowdisplaybreaks
\setlist[enumerate]{itemsep=0.45em, topsep=0.35em, leftmargin=2.4em}
\setlist[itemize]{itemsep=0.25em, topsep=0.25em, leftmargin=1.9em}

\definecolor{mainblue}{RGB}{36,83,140}
\definecolor{lightblue}{RGB}{241,246,252}
\definecolor{accent}{RGB}{153,76,0}
\definecolor{lightgray}{RGB}{246,246,246}
\definecolor{rulegray}{RGB}{110,118,128}

\titleformat{\section}
  {\Large\bfseries\color{mainblue}}
  {\thesection.}
  {0.45em}
  {}
  [\vspace{0.2em}\color{rulegray}\titlerule]
\titlespacing*{\section}{0pt}{1.1em}{0.7em}

\pagestyle{fancy}
\fancyhf{}
\fancyhead[L]{\small 染色方法与棋盘问题周测 B 卷}
\fancyhead[C]{\small 组合专题：多色染色与不变量}
\fancyhead[R]{\small 刘文博}
\fancyfoot[C]{\small 第 \thepage 页 / 共 \pageref{LastPage} 页}
\renewcommand{\headrulewidth}{0.55pt}
\renewcommand{\footrulewidth}{0.35pt}

\newtcolorbox{infobox}[1]{
  breakable,
  colback=lightblue,
  colframe=mainblue,
  boxrule=0.7pt,
  arc=1mm,
  title=\textbf{#1},
  fonttitle=\bfseries
}

\newenvironment{solution}{\par\noindent\textbf{参考解答：}}{\hfill$\square$\par}
\newcommand{\answerspace}{\vspace{2.3cm}}
\newcommand{\biganswerspace}{\vspace{3.1cm}}

\begin{document}

\thispagestyle{empty}
\begin{center}
{\fontsize{22}{28}\selectfont\bfseries 染色方法与棋盘问题周测 B 卷\par}
\vspace{0.4cm}
{\large 适合小学高年级至初中低年级数学竞赛训练\par}
\end{center}

\begin{infobox}{考试说明}
\begin{itemize}
  \item 测试时间：75--80 分钟
  \item 满分：100 分
  \item B 卷侧重“必要条件不是充分条件”“多色染色”和“奇偶性不变量”，比 A 卷更强调思路判断。
\end{itemize}
\end{infobox}

\section{试题}

\begin{enumerate}
  \item （12 分）请给出一个“黑格数与白格数相等，但仍然不能被 $1\times 2$ 骨牌完全覆盖”的图形例子，并说明原因。
  \answerspace

  \item （12 分）把 $6\times 6$ 棋盘的左上角和右下角去掉，判断剩下的图形能否被 $1\times 2$ 骨牌完全覆盖，并说明理由。
  \answerspace

  \item （12 分）在 $7\times 7$ 棋盘上，若每一步都只能走到上下左右相邻的格子，并且每个格子恰好经过一次，证明这条路径的起点和终点一定同色。
  \answerspace

  \item （12 分）国际象棋中的马能否在走了 $17$ 步后回到与出发点同色的格子？请说明理由。
  \answerspace

  \item （14 分）判断并说明：$6\times 6$ 棋盘能否被 $12$ 个 $1\times 3$ 长条完全覆盖？
  \biganswerspace

  \item （14 分）请举出一个“格子总数是 $3$ 的倍数，但仍然不能被 $1\times 3$ 长条完全覆盖”的例子，并说明理由。
  \biganswerspace

  \item （12 分）一个 $8\times 8$ 棋盘最开始只有 $1$ 个黑格，其余 $63$ 个格子都是白格。每次操作允许任选两个相邻格子，把这两个格子的颜色同时翻转。问经过若干次操作后，能否把整个棋盘都变成白色？
  \answerspace

  \item （12 分）请说明为什么“颜色数量刚好平衡”往往只是必要条件而不是充分条件，并各举一个骨牌题或长条覆盖题中的例子支持你的说明。
  \answerspace
\end{enumerate}

\begin{center}
\begin{tikzpicture}[scale=0.58]
  \node[draw=none,fill=none,font=\bfseries\color{mainblue}] at (6.2,8.8) {附图参考};
  \foreach \x in {0,...,5} {
    \foreach \y in {0,...,5} {
      \pgfmathtruncatemacro{\c}{mod(\x+\y,2)}
      \ifnum\c=0 \fill[black!18] (\x,\y) rectangle ++(1,1); \else \fill[white] (\x,\y) rectangle ++(1,1); \fi
      \draw[black!55] (\x,\y) rectangle ++(1,1);
    }
  }
  \draw[accent,line width=1.1pt] (0,5) rectangle ++(1,1);
  \draw[accent,line width=1.1pt] (5,0) rectangle ++(1,1);
  \draw[accent,line width=1.0pt] (0,5)--(1,6) (1,5)--(0,6) (5,0)--(6,1) (6,0)--(5,1);

  \begin{scope}[xshift=8.3cm,scale=0.95]
    \foreach \x in {0,...,4} {
      \foreach \y in {0,...,4} {
        \pgfmathtruncatemacro{\c}{mod(\x+\y,2)}
        \ifnum\c=0 \fill[black!18] (\x,\y) rectangle ++(1,1); \else \fill[white] (\x,\y) rectangle ++(1,1); \fi
        \draw[black!55] (\x,\y) rectangle ++(1,1);
      }
    }
    \fill[black!70] (1.5,1.5) circle (0.16);
    \draw[green!45!black,line width=1.2pt,-{Stealth[length=2.6mm]}] (1.65,1.65) to[out=25,in=210] (3.35,2.35);
    \fill[white] (3.5,2.5) circle (0.16);
  \end{scope}
\end{tikzpicture}
\end{center}

\newpage
\section{详细解答}

\begin{enumerate}
  \item \begin{solution}
  可以取 $3\times 3$ 棋盘的四个角格组成的图形。

  这个图形中有 $2$ 个黑格、$2$ 个白格，黑白数量相等。

  但是四个角格彼此都不相邻，所以在这个图形内部根本放不下一块 $1\times 2$ 骨牌。既然一块都放不下，就不可能完全覆盖。

  这说明：黑白格数相等只是必要条件，不是充分条件。
  \end{solution}

  \item \begin{solution}
  不能覆盖。

  左上角和右下角同色，因此把这两个角去掉以后，黑白格数不再相等。

  原来 $6\times 6$ 棋盘中有 $18$ 个黑格、$18$ 个白格。去掉两个同色角格后，变成
  \[
  16\text{ 个某种颜色，}18\text{ 个另一种颜色。}
  \]

  但每块 $1\times 2$ 骨牌都必须覆盖一黑一白，所以若能完全覆盖，就应该黑白数量相等。

  矛盾，因此不能覆盖。
  \end{solution}

  \item \begin{solution}
  $7\times 7$ 棋盘共有
  \[
  7\times 7=49
  \]
  个格子，是奇数个。

  一条满足条件的路径要把这 $49$ 个格子都经过一次，因此它经过的格子总数也是奇数个。

  又因为沿格边前进时颜色黑白交替，所以：
  \begin{itemize}
    \item 经过偶数个格子时，起点终点异色；
    \item 经过奇数个格子时，起点终点同色。
  \end{itemize}

  现在经过的是 $49$ 个格子，属于奇数情形，因此起点和终点一定同色。
  \end{solution}

  \item \begin{solution}
  不能。

  马每走一步都会改变颜色：从黑格跳到白格，或从白格跳到黑格。

  因而：
  \begin{itemize}
    \item 走奇数步后，一定落在与起点异色的格子上；
    \item 走偶数步后，才会落在与起点同色的格子上。
  \end{itemize}

  因为 $17$ 是奇数，所以马走了 $17$ 步后，不可能回到与出发点同色的格子。
  \end{solution}

  \item \begin{solution}
  能覆盖。

  因为 $6\times 6$ 棋盘每一行都有 $6$ 个格子，而 $6$ 正好可以分成两个长度为 $3$ 的小段。所以每一行都能用两个横放的 $1\times 3$ 长条盖满。

  六行一共需要
  \[
  6\times 2=12
  \]
  个 $1\times 3$ 长条。

  因此 $6\times 6$ 棋盘可以被 $12$ 个 $1\times 3$ 长条完全覆盖。
  \end{solution}

  \item \begin{solution}
  一个经典例子是：$8\times 8$ 棋盘去掉左上角以后，剩下 $63$ 个格子。

  首先，$63$ 的确是 $3$ 的倍数，这说明“格子总数”这个条件没有问题。

  但如果把格子按 $(i+j)\bmod 3$ 染成三种颜色，那么完整 $8\times 8$ 棋盘中三种颜色个数分别为
  \[
  21,\ 22,\ 21.
  \]

  去掉左上角后变成
  \[
  20,\ 22,\ 21.
  \]

  任意一个横放或竖放的 $1\times 3$ 长条都覆盖三种颜色各一个，所以若能完全覆盖，则三种颜色总数必须相等。

  现在三种颜色个数不相等，因此不能覆盖。

  这就给出了一个“格子总数是 $3$ 的倍数，但仍不能覆盖”的例子。
  \end{solution}

  \item \begin{solution}
  不能。

  关注黑格总数的奇偶性。开始时只有 $1$ 个黑格，所以黑格数是奇数。

  每次翻转两个相邻格时：
  \begin{itemize}
    \item 如果原来两格都是白色，黑格数增加 $2$；
    \item 如果原来两格都是黑色，黑格数减少 $2$；
    \item 如果原来一黑一白，黑格数不变。
  \end{itemize}

  无论哪种情况，黑格数的变化都是 $0$ 或 $\pm 2$，所以黑格数的奇偶性保持不变。

  既然一开始黑格数是奇数，那么以后任何时刻黑格数都还是奇数。

  但整个棋盘都变成白色时，黑格数是 $0$，为偶数。矛盾。

  因此不可能把整个棋盘都变成白色。
  \end{solution}

  \item \begin{solution}
  因为“颜色数量平衡”只说明每一步在颜色数量上\textbf{没有立刻矛盾}，并没有保证图形的相邻关系、连通情况和局部结构一定允许覆盖。

  例如骨牌题中：
  \begin{itemize}
    \item $3\times 3$ 棋盘四个角格组成的图形，有 $2$ 黑 $2$ 白，颜色数量平衡；
    \item 但四个角彼此都不相邻，一块骨牌也放不进去，所以仍然无法覆盖。
  \end{itemize}

  再例如 $1\times 3$ 长条题中：
  \begin{itemize}
    \item 在某些图形里，三种颜色个数也许相等；
    \item 但若图形被切得很零散，或局部长度根本放不下一条 $1\times 3$，仍然可能无法覆盖。
  \end{itemize}

  因此，颜色数量平衡通常只是必要条件。想得到“可以覆盖”的结论，往往还要进一步给出具体构造或更强的结构分析。
  \end{solution}
\end{enumerate}

\end{document}
